\chapter*{Natureza e controle deste documento}
\addcontentsline{toc}{chapter}{Natureza e controle deste documento}

\section*{O que este documento é}

Este volume descreve como o \NomeDoProduto\ foi construído e verificado: a arquitetura, as regras de codificação que o projeto impõe, o processo de trabalho, o uso de inteligência artificial no desenvolvimento, o catálogo de testes automatizados, os resultados do plano de testes e as avaliações de experiência de uso. É um documento interno de engenharia e, ao mesmo tempo, um guia: cada capítulo diz o que foi feito, por que, como foi aplicado neste projeto, qual a evidência e que lição ficou.

Os números citados vêm do repositório, dos registros de execução dos testes e do histórico do controle de versão, sem arredondamento. O que não foi executado está dito como não executado.

\section*{Controle do documento}

\begin{table}[H]
\tabelapequena
\begin{tabularx}{\textwidth}{L{4.1cm}Y}
\toprule
\linhacab \cab{Item} & \cab{Informação}\\
\midrule
Versão e data & Versão 1.0, de \DataBase\\
Elaboração & \NomeDoAutor, \CargoDoAutor\\
Destinatários & \NomeDoDestinatario\\
Classificação & Uso interno\\
Base de código & Repositório \cod{nome-do-repositorio}, revisão \cod{abc1234}, de DATA\\
Período dos testes & DATA-INICIAL a DATA-FINAL\\
Documento relacionado & Volume~I, \textit{Tratamento de Dados Pessoais no \NomeDoProduto}\\
\bottomrule
\end{tabularx}
\end{table}

\section*{Como ler}

A Parte~I explica o produto por dentro; a Parte~II, como o trabalho é feito; a Parte~III, os testes e o que provaram; a Parte~IV, as avaliações de uso. O Capítulo~\ref{cap:guia} reúne as práticas em forma de guia, com uma lista de verificação para começar um projeto novo. Os nomes completos de todos os testes estão no Anexo~\ref{anexo:testes}.

Duas caixas aparecem ao longo do texto: \enquote{Como aplicar} descreve uma prática de forma reutilizável, e \enquote{Lição aprendida} registra o que este projeto aprendeu ao aplicá-la ou ao deixar de aplicá-la.
